\section{Projekt rozwiązania i implementacja}
\subsection{Kod genetyczny osobnika}
W niniejszym dokumencie kod genetyczny osobnika będzie reprezentował
architekturę sieci i jej parametry uczenia.
Kod genetyczny będzie składał się z trzech części:
\begin{itemize}
  \item Część binarna:
    \begin{itemize}
      \item liczba warstw ukrytych: 5 bitów $\left[1, 32\right]$
      \item ziarno generatora liczby neuronów: 10 bitów $\left[0, 1023\right]$
      \item ziarno generatora typów aktywacji: 10 bitów $\left[0, 1023\right]$
      \item optymalizator: 1 bit $\left[0, 1\right]$ (0 \textrightarrow SGD, 1 \textrightarrow Adam)
      \item wskaźnik uczenia się (learning rate): 12 bitów $\left[2^{-13}, 0.5\right]$
      % TODO: maybe this version [2^(-13) - 0.5] → [2^(-13) - 1.]
      % \item wskaźnik uczenia się (learning rate): 13 bitów $\left[2^{-13}, 1.0)\right]$
      \item liczba epok: 7 bitów $\left[1, 128\right]$
      \item wielkość partii: 10 bitów $\left[1, 1024\right]$
    \end{itemize}
  \item Część list - generowane po części binarnej;
    długości list są nie mniejsze niż liczba warstw ukrytych,
    przy czym długości wygenerowanych są równe tej liczbie;
    minimalne i maksymalne długości list są zależne
    od minimalnych i maksymalnych wartości liczby warstw ukrytych:
    \begin{itemize}
      \item lista ilości neuronów w warstwach ukrytych -\\
        \hspace*{1em}liczba elementów - taka sama:
          $\left[1, 128\right]$,\\
        \hspace*{1em}zawiera liczby z zakresu:
          $\left[2, 724\right]$,\\
        % ze wzgędu na możliwe ograniczenia pamięci dla sieci: .5GB
        % .5÷128÷8×1024×1024×1024=524288 -> √524288 ~ 724
        \hspace*{1em}każdy element składa się z 10 bitów.
      \item lista typów aktywacji w warstwach ukrytych -\\
        \hspace*{1em}liczba elementów - o 1 większa
          (uwzględnia warstwę wyjściową):
          $\left[2, 129\right]$,\\
        \hspace*{1em}zawiera liczby z zakresu:
          $\left[0, 3\right]$
          (%
            0 \textrightarrow relu,
            1 \textrightarrow tanh,
            2 \textrightarrow sigmoid,
            3 \textrightarrow leaky relu%
          ),\\
        \hspace*{1em}każdy element składa się z 2 bitów.
    \end{itemize}
\end{itemize}
Ważnym mechanizmem podczas czytania kodu genetycznego
w celu utworzenia fenotypu są ziarna generatora.

\noindent\\
Jeżeli liczba warstw ukrytych z części binarnej,
wynosiłoby więcej niż liczba elementów w listach,
to~należy uzupełnić listy ilości neuronów
i typów aktywacji do danej liczby elementów do czego zostanie
użyty generator z odpowiednim ziarnem.
Następnie kod genetyczny zostanie uaktualniony o nowe elementy list.

\noindent\\
Jeżeli liczba warstw ukrytych jest mniejsza niż liczba elementów w listach,
to nadmiarowe elementy zostaną zignorowane w fenotypie.

\subsubsection{Selekcja}
Selekcja będzie odbywała się metodą ruletki proporcjonalnej do przystosowania.
W przypadku, gdyby wszystkie osobniki miały przystosowanie równe 0,
to zostanie użyta selekcja losowa.

% Krzyżowanie będzie odbywało się różnie dla części binarnej
% i listowej kodu genetycznego.
\subsubsection{Krzyżowanie - część binarna}
W przypadku części binarnej zostanie użyte krzyżowanie jednopunktowe,
gdzie losowo zostanie wybrany punkt podziału i
części kodu genetycznego zostaną zamienione między dwoma rodzicami.

\noindent\\
Matematyczny opis krzyżowania, gdzie\\
$n$ - długość genotypu (liczba bitów),\\
$R\left(p, k\right)$ - funkcja losowa o rozkładzie jednostajnym dyskretnym
  na zbiorze $\mathbb{C} \cap \left[p, k\right]$,\\
$r_1\left[i\right]$ i $r_2\left[i\right]$ -
  rodzice 1 i 2 dla
  $i \in \mathbb{N} \cap \left[1, n\right]$,\\
$d_1\left[i\right]$ i $d_2\left[i\right]$ -
  dzieci 1 i 2 dla
  $i \in \mathbb{N} \cap \left[1, n\right]$:

\noindent\\
Wylosowanie punktu przecięcia i końcowa postać kodu dzieci:\\
$\displaystyle
  p_r=R\left(0, n\right)\\
  d_1\left[i\right]=\begin{cases}
    r_1[i]
      &, i \le p_r\\
    r_2[i]
      &, i>p_r\\
  \end{cases}\\
  d_2\left[i\right]=\begin{cases}
    r_2[i]
      &, i \le p_r\\
    r_1[i]
      &, i>p_r\\
  \end{cases}
$

\subsubsection*{Elitaryzm}
Zrezygnowano z elitaryzmu w celu zwiększenia
różnorodności populacji i ograniczenia
ryzyka przedwczesnej zbieżności.

\subsubsection{Krzyżowanie - część listowa}
W przypadku części listowej zostanie użyte krzyżowanie zmiennej długości,
gdzie losowo są wybrane punkty podziału w obu listach\,\sscite{synapsing_varlen_crossover}.
Natomiast jest on inny od opisanego w przypisie algorytmu,
ponieważ o ile pierwszy punkt podziału jest wybierany z pełnego zakresu
(biorąc pod uwagę wszystkie bity listy),
to punkt podziału drugiej listy jest wybierany na podstawie
wylosowanego elementu listy, by zachować spójność długości list.
Dodatkowo algorytm pilnuje by lista nie miała mniej elementów,
niż minimalna liczba, a regulowanie do maksymalnej wartości
jest wykonywane przez przycinanie listy po krzyżowaniu.

\noindent\\
Oznacza to, że pierwsze losowanie,
poza decydowaniem konkretnego przecięcia, decyduje o tym czy
i w którym punkcie nastąpi podział elementów w drugiej liście.

\noindent\\
Matematyczny opis krzyżowania, gdzie\\
$e_s$ - wielkość pojedyńczego elementu,\\
$l_{\min}$ i $l_{\max}$ - minimalna i maksymalna liczba elementów,\\
$R\left(p, k\right)$ - funkcja losowa o rozkładzie jednostajnym dyskretnym
  na zbiorze $\mathbb{C} \cap \left[p, k\right]$,\\
$n_{r_1}$ i $n_{r_2}$ - liczba elementów dla rodzica 1 i 2,\\
$r_1\left[i_{r_1}\right]$ i $r_2\left[i_{r_2}\right]$ -
  rodzice 1 i 2 odpowiednio dla
  $i_{r_1} \in \mathbb{N} \cap \left[1, e_s n_{r_1}\right]$
  i $i_{r_2} \in \mathbb{N} \cap \left[1, e_s n_{r_2}\right]$,\\
$n_{d_1}$ i $n_{d_2}$ - liczba elementów dla dziecka 1 i 2,\\
$d_1\left[i_{d_1}\right]$ i $d_2\left[i_{d_2}\right]$ -
  dzieci 1 i 2 odpowiednio dla
  $i_{d_1} \in \mathbb{N} \cap \left[1, e_s n_{d_1}\right]$
  i $i_{d_2} \in \mathbb{N} \cap \left[1, e_s n_{d_2}\right]$:

\noindent\\
Losowanie punktu przecięcia dla pierwszego rodzica i obliczenia dla niego:\\
$\displaystyle
  p_{r_1}=R\left(0, e_s n_{r_1}\right)\\
  l_b=p_{r_1} \mod e_s\\
  l_{p_b}=\begin{cases}
    1 &, l_b>0\\
    0 &, l_b=0
  \end{cases}\\
  l_{p_{d_1}}=\left\lfloor
    \frac{p_{r_1}}{e_s}
  \right\rfloor + l_{p_b}\\
  l_{k_{d_2}}=n_{r_1}-\left\lfloor
    \frac{p_{r_1}}{e_s}
  \right\rfloor
$

\noindent\\
Przygotowanie do losowania i losowanie punktu przecięcia dla drugiego rodzica:\\
$\displaystyle
  p_{r_{2_p}}=\begin{cases}
    0
      &, l_{k_{d_2}} \ge l_{\min}\\
    l_{\min}-l_{k_{d_2}}
      &, l_{k_{d_2}}<l_{\min}\\
  \end{cases}\\
  p_{r_{2_k}}=\begin{cases}
    n_{r_2}-l_{p_b}
      &, l_{p_{d_1}} \ge l_{\min}\\
    n_{r_2}-\left(l_{\min}-l_{p_{d_1}}+l_{p_b}\right)
      &, l_{p_{d_1}}<l_{\min}\\
  \end{cases}\\
  p_{r_2}=e_s R\left(p_{r_{2_p}}, p_{r_{2_k}}\right) + l_b
$

\noindent\\
Końcowa długość i postać kodu dzieci:\\
$\displaystyle
  {\Delta p_r}=p_{r_1}-p_{r_2}\\
  n_{d_1}=\begin{cases}
    \frac{p_{r_1}+\left(e_s n_{p_2}-p_{r_2}\right)}{e_s}=
      n_{p_2}+\frac{\left(p_{r_1}-p_{r_2}\right)}{e_s}=
      n_{p_2}+\frac{{\Delta p_r}}{e_s}
        &, n_{p_2}+\frac{{\Delta p_r}}{e_s} \le l_{\max}\\
    l_{\max}
      &, n_{p_2}+\frac{{\Delta p_r}}{e_s}>l_{\max}\\
  \end{cases}\\
  n_{d_2}=\begin{cases}
    \frac{p_{r_2}+\left(e_s n_{p_1}-p_{r_1}\right)}{e_s}=
      n_{p_1}-\frac{\left(p_{r_1}-p_{r_2}\right)}{e_s}=
      n_{p_1}-\frac{{\Delta p_r}}{e_s}
        &, n_{p_1}-\frac{{\Delta p_r}}{e_s} \le l_{\max}\\
    l_{\max}
      &, n_{p_1}-\frac{{\Delta p_r}}{e_s}>l_{\max}\\
  \end{cases}\\
  d_1\left[i_{d_1}\right]=\begin{cases}
    p_1[i_{d_1}]
      &, i_{d_1} \le p_{r_1}\\
    p_2[i_{d_1}-p_{r_1}+p_{r_2}]=
      p_2[i_{d_1}-{\Delta p_r}]
      &, i_{d_1}>p_{r_1}\\
  \end{cases}\\
  d_2\left[i_{d_2}\right]=\begin{cases}
    p_2[i_{d_2}]
      &, i_{d_2} \le p_{r_2}\\
    p_1[i_{d_2}-p_{r_2}+p_{r_1}]=
      p_1[i_{d_2}+{\Delta p_r}]
      &, i_{d_2}>p_{r_2}\\
  \end{cases}
$

\subsubsection{Mutacja}
Mutacja będzie odbywała się przez odwracanie losowych bitów.
Decyzja o mutacji będzie podejmowana dla każdego bitu
z osobna z ustalonym prawdopodobieństwem mutacji.

\subsubsection{Ocena przystosowania}
Ocena przystosowania osobnika będzie odbywała się poprzez nauczenie
pięciu sieci o parametrach i z parametrami wynikającymi
z fenotypu osobnika.
Obliczone dokładności na danych testowych zostaną uśrednione
i przekształcone na przystosowanie za pomocą wzoru:
$$
\mbox{fit} = \frac{\mbox{dokładność}}{1 - \mbox{dokładność}} % TODO: change it
$$

\noindent\\
Jeżeli użytkownik poda tylko dane treningowe i testowe, to:\\
Uczenie będzie odbywało się na tylko danych treningowych.
Po zakończeniu uczenia zostanie obliczona dokładność na danych testowych.
Zwrócona sieć będzie tą z ostatniej epoki uczenia.

\noindent\\
Jeżeli użytkownik poda dodatkowo dane walidacyjne, to:\\
Uczenie będzie odbywało się na danych treningowych,
a po każdej epoce będzie sprawdzana dokładność na danych walidacyjnych.
Zwrócona sieć będzie tą z epoki,
która dała najlepszą dokładność na danych walidacyjnych.
Jeżeli dokładność ta będzie reprezentowana przez kilka epok,
to zostanie wybrana ta z najwcześniejszej epoki. % może lepsze z najpóźniejszej?

% \begin{minted}[frame=single]{python3}
%   gen=(
%     '000000010000010010110001010100000010011001010000110010',
%     [64,32,8,16,4],
%     [0,1,0,2,1]
%   )
% \end{minted}
% Powyższy kod genetyczny reprezentuje następujący fenotyp:
% \begin{minted}[frame=single]{python3}
%   # fenotyp=(
%   #   ilość ukrytych warstw: 1-1024,
%   #   ziarno generatora liczby neuronów: 0-1023,
%   #   ziarno generatora typów aktywacji: 0-1023,
%   #   optymalizator: 0->SGD, 1->Adam,
%   #   wskaźnik uczenia się (learning rate): 1/2^13-1.0 [13bit],
%   # ? wskaźnik uczenia się (learning rate): 1/2^13-0.5 [12bit],
%   #   liczba epok: 1-1024,
%   #   ilość neuronów w warstwach: lista[2-1 048 577],
%   #   typy aktywacji w warstwach: lista[0-2] (0->relu, 1->tanh, 2->sigmoid),
%   # )
%   # gen=(
%   #   '|0000000100|0001001011|0001010100|0|0001001100101|0000110001|',
%   # #  |   4+1    |    75    |    84    |0|    613+1    |   49+1   |
%   # #  |    5     |    75    |    84    |0|  614/2^13   |    50    |
%   # #  |    5     |    75    |    84    |0| 0.074951172 |    50    |
%   #   [64,32,8,16,4],
%   #   [0,1,0,2,1]
%   # )
%   fenotyp=(
%     5,
%     75,
%     84,
%     0,
%     0.074951172,
%     50,
%     [64, 32, 8, 16, 4],
%     [0, 1, 0, 2, 1],
%   )
% \end{minted}